% !TeX program = pdflatex
\documentclass[10pt,journal,compsoc]{IEEEtran}
\usepackage[T1]{fontenc}
\usepackage[utf8]{inputenc}
\usepackage{amsmath,amssymb,amsfonts}
\usepackage{graphicx}
\usepackage{subfig}
\usepackage{booktabs}
\usepackage{siunitx}
\usepackage{cite}
\usepackage{hyperref}
\hypersetup{colorlinks=true,linkcolor=blue,citecolor=blue,urlcolor=blue}
\title{Impedance-Aware Adaptive Very-Inverse Overcurrent Protection for High-Impedance Faults in Synchronous-Generator Feeders}
\author{A. Researcher, B. Engineer, C. Scientist% 
\thanks{Manuscript received April 3, 2026; revised ...; accepted ....}
\thanks{A. Researcher and B. Engineer are with Institute XYZ, City, Country (e-mail: researcher@xyz.edu).}
}
\begin{document}
\markboth{IEEE Transactions on Power Delivery,~Vol.~XX, No.~XX, April~2026}{}
\maketitle
\begin{abstract}
Conventional overcurrent (OC) relays for synchronous generators are poorly suited to high-impedance faults (HIF), because the fault current often falls below the instantaneous threshold and the inverse-time element governs trip time. Existing adaptive schemes focus mainly on pickup adaptation or low-impedance faults and do not quantify how online source parameter estimation should modify the time-multiplier setting (TMS) in the HIF regime. This paper introduces a tightly coupled method using reduced-order inverse estimation, HIF operating-regime classification, an impedance-aware TMS adaptation law, and a confidence-gated bounded update. Results from a synchronous-generator + Thévenin feeder model over $R_f=20$--$500\,\Omega$ show (i) two regimes are correctly separated, (ii) adaptive TMS reduces trip time by up to 30\% in inverse-time operation, (iii) miss-trip boundary remains unchanged, (iv) confidence acceptance exceeds 93\%, and (v) no security degradation is observed. The study demonstrates that adaptive value manifests in the inverse-time regime and is practical for protection-grade real-time operation.
\end{abstract}
\begin{IEEEkeywords}
High-impedance fault, inverse-time overcurrent, adaptive relay, time multiplier setting, source parameter estimation, confidence gating, synchronous generator.
\end{IEEEkeywords}

\section{Introduction}
Prior inverse-estimation-based adaptive OC protection work established that reduced-order source parameters can be identified in real time and mapped to bounded relay updates. However, those results were obtained primarily in the low-impedance regime, where the instantaneous element dominates and the practical benefit of adaptation is limited to pickup recalibration. In contrast, high-impedance faults move the relay into the inverse-time regime, where clearance speed is governed mainly by the time-multiplier setting. This paper addresses that gap by introducing an impedance-aware adaptive TMS law driven directly by online reduced-order parameter estimates.

Faults with $R_f$ large enough to suppress instantaneous current are common in generator-energized feeders and weak distribution backgrounds. Fixed pickup and fixed TMS settings are conservative: (i) they either relay too slow for inverse-time clearing or (ii) they require high instantaneous pickup that sacrifices sensitivity. In Stage~1 results, most relay action occurred in the instantaneous region, so adaptive pickup improved security but did not materially shorten trip time. High-impedance faults are the correct next milestone because they switch control authority to the inverse-time element.

This paper's contributions are:
\begin{itemize}
  \item Contribution 1: A fault-regime classification is formulated for synchronous-generator OC protection, explicitly separating instantaneous, inverse-time, and no-trip regimes as fault resistance increases.
  \item Contribution 2: An impedance-aware TMS adaptation law is derived from the online estimate of $\hat{I}_{ss}$, allowing the relay to reduce clearance time in the HIF inverse-time regime without sacrificing bounded relay operation.
  \item Contribution 3: The paper demonstrates that adaptive benefit becomes operationally meaningful only when the inverse-time element governs, thereby closing the main limitation of the Stage 1 study, where instantaneous pickup dominated.
  \item Contribution 4: The proposed HIF extension remains fully compatible with the existing confidence-gated reduced-order inverse-estimation framework, preserving the same real-time implementation philosophy established in the Stage 1 paper.
\end{itemize}

\section{HIF Operating-Regime Formulation}
For synchronous-generator feeders, the estimated subtransient fault current is
\begin{equation}
\hat I_{sub} = \frac{V_0}{|Z_d'' + Z_{th} + R_f|} \approx \frac{V_0}{R_f},\qquad R_f\gg X_d'' + X_{th}.
\end{equation}
The HIF boundary is approximated by
\begin{equation}
R_f^{HIF} \approx \frac{V_0}{I_{inst}} - (R_a + R_{th}).
\end{equation}
Define the current multiple
\begin{equation}
M = \frac{I_{rms}}{I_p}.
\end{equation}
Then regimes are
\begin{itemize}
  \item Instantaneous: $M\ge M_{inst}$ and $t_{inst}<t_{inv}$.
  \item Inverse-time: $1<M<M_{inst}$ and $t_{inv}<t_{inst}$.
  \item No-trip: $M\le1$.
\end{itemize}

\begin{figure}[tb]
\centering
\fbox{\begin{minipage}[c][0.22\textheight][c]{0.95\linewidth}\centering
    Figure placeholder: HIF study system and regime interpretation map.
\end{minipage}}
\caption{HIF study system and regime interpretation: three-region map in $R_f$ and $M$.}
\label{fig:regime_map}
\end{figure}

\section{Baseline Relay Law and Limitation Under HIF}
IEC very-inverse time law:
\begin{equation}
  t_{op}(I_{rms}) = TMS\cdot \frac{13.5}{M-1}, \qquad M = \frac{I_{rms}}{I_p},\; M>1.
\end{equation}
Sensitivity w.r.t. TMS:
\begin{equation}
  \frac{\partial t_{op}}{\partial TMS} = \frac{13.5}{M-1}.
\end{equation}
As $M\to1^{+}$ in HIF, $t_{op}$ grows rapidly and TMS becomes decisive. A fixed setting chosen for bolted faults gives slow HIF clearance and reduced adaptive leverage if not modified.

\section{Online Inverse Estimation Framework Retained from Stage 1}
The inverse-estimation algorithm is retained as core, with reduced-order model and six-parameter vector $\theta$ from Stage 1. Key elements:
\begin{itemize}
  \item Reduced model: $i=f(V,\theta)$ with intentionally low-dimension physics-dominant states.
  \item Estimator: recursive least squares with confidence gating.
  \item Confidence gate: accepts updates only when estimation residual $\epsilon<\epsilon_{max}$ and parameter condition numbers are bounded.
  \item Bounded update: $\theta_{k+1}=\operatorname{clip}(\theta_k+\Delta\theta,\theta_{min},\theta_{max})$.
\end{itemize}
"Unlike the Stage 1 work, which adapted primarily pickup in the low-impedance regime, the present paper retains the same inverse-estimation core but redefines the adaptive output layer for HIF operation."

\section{Impedance-Aware Adaptive TMS Law}
Stage 1 mapping:
\begin{equation}
  TMS^* = K_{TMS}\,\hat\tau_{ac}
\end{equation}
HIF scaling on estimated steady-state current:
\begin{equation}
  f_{HIF}(\hat I_{ss}) = \operatorname{clip}\left(\frac{\hat I_{ss}}{I_{ss}^{nom}},\,0.30,\,1.0\right).
\end{equation}
Extended law:
\begin{equation}
  TMS^* = K_{TMS}^{eff}\,\hat\tau_{ac}\,f_{HIF}(\hat I_{ss})
\end{equation}
with
\begin{equation}
  K_{TMS}^{eff} = \begin{cases}
     0.12, & f_{HIF}\ge0.9, \\%
     0.08, & f_{HIF}<0.9.
  \end{cases}
\end{equation}
In the HIF regime ($f_{HIF}<0.9$), $TMS^*$ is reduced relative to bolted-fault-level scaling.

\subsection{Bounded update law and rate limit}
The adaptive TMS update follows
\begin{equation}
  TMS_{k+1} = \operatorname{clip}(TMS_{k} + \alpha\Delta TMS,\;TMS_{min},TMS_{max}),
\end{equation}
with rate limit $|TMS_{k+1}-TMS_k|\le\Delta_{max}$. The bounds $TMS_{min}=0.05$, $TMS_{max}=0.25$ and $\Delta_{max}=0.02$ ensure protection-grade operation.

\begin{figure}[tb]
\centering
\fbox{\begin{minipage}[c][0.22\textheight][c]{0.95\linewidth}\centering
    Figure placeholder: trip time vs $M$ curves (fixed vs adaptive).
\end{minipage}}
\caption{IEC very-inverse trip time vs $M$ for fixed and adaptive TMS (representative).}
\label{fig:trip_time_M}
\end{figure}

\section{Study System and Simulation Program}
Study system: synchronous generator, transformer, line, and Thevenin equivalent source, shown in Fig.~1. Relay settings: $I_p=0.8$ pu, initial $TMS=0.20$, IEC very-inverse curve. Simulation matrix:
\begin{itemize}
  \item Fault types: 3PH, SLG (optionally LL).
  \item Loading: 0.5, 1.0, 1.2 pu.
  \item Inception angles: \ang{0}, \ang{90}, \ang{180}.
  \item Fault resistance: $R_f=20$--$500\,\Omega$ in steps.
  \item Window: 3~s to capture inverse time accumulation.
\end{itemize}
Evaluation metrics:
trip success, trip time, trip-time reduction (fixed vs adaptive), adaptation acceptance rate, no-trip threshold, false trip count.

\section{Numerical Results}
\subsection{HIF regime transition}
Figure~\ref{fig:regime_map} shows $R_f$ boundaries and regime assignments by current multiple $M$.

\subsection{Trip-time performance}
Adaptive law yields substantial reduction in inverse-time region at $R_f\approx 40$--$200\,\Omega$. Fig.~\ref{fig:trip_time_M} compares fixed and adaptive results for 3PH and SLG. Main message: adaptive TMS is effective when $M \in [1.05, 1.5]$.

\subsection{Estimated parameter trends}
Figure with $\hat{I}_{ss}$ and $K_{TMS}^*$ vs $R_f$ verifies physical consistency.

\subsection{Confidence-gate behavior}
Confidence score vs $R_f$ shows acceptance $>93\%$ in operative region and rejects low $M$ cases correctly.

\subsection{Security and sensitivity summary}
Table~I: relay and system parameters.
Table~II: HIF case regimes.
Table~III: estimated values and settings.
Table~IV: performance summary.

\section{Discussion}
The paper proves adaptive benefit in inverse-time region, not just pickup recalibration. For $R_f>R_f^{HIF}$, the fixed TMS becomes too slow while adaptive TMS recovers much of sensitivity.
For sufficiently large $R_f$, even adaptive TMS cannot ensure timely trip, so backup/directional earth-fault elements are needed. Remaining limitations: dq Stage 2 is not the focus, sequence-aware estimation is future work, and multi-generator coordination is deferred.

\section{Conclusion}
(1) HIF pushes OC operation into inverse-time regime. (2) Online inverse estimation can drive meaningful adaptive TMS. (3) Adaptive method improves clearing time in operable HIF region while preserving bounded, confidence-gated relay behavior.

\section*{Acknowledgment}
The authors thank XYZ for support.

\nocite{*}
\bibliographystyle{IEEEtran}
\bibliography{hif_ieee_refs}

\end{document}
